% Part:sets-functions-relations
% Chapter: size-of-sets
% Section: non-enumerability

\documentclass[../../../include/open-logic-section]{subfiles}

\begin{document}

\olfileid[pt]{sfr}{siz}{nen}

\olsection{Conjuntos \printtoken{P}{nonenumerable}}

\begin{editorial}
  Esta seção prova a não enumerabilidade de $\Bin^\omega$ e
  $\Pow{\PosInt}$ usando a definição na \olref[enm]{sec}. Foi
  concebida para ser um pouco mais elementar e um pouco mais detalhada
  que a versão na \olref[enm-alt]{sec}
\end{editorial}

Alguns conjuntos, como o conjunto $\PosInt$ dos inteiros positivos, são
infinitos. Até agora vimos exemplos de conjuntos infinitos que eram todos
!!{enumerable}. Contudo, há também conjuntos infinitos que não
têm essa propriedade. Tais conjuntos são chamados de \emph{!!{nonenumerable}}.

Em primeiro lugar, talvez já seja surpreendente que existam
conjuntos !!{nonenumerable}. Para qualquer conjunto !!{enumerable}~$A$ há
uma função !!{surjective} $f \colon \PosInt \to A$. Se um conjunto é
!!{nonenumerable}, não há tal função. Isto é, nenhuma função
mapeando os infinitos !!{element}s de~$\PosInt$ para~$A$ pode
esgotar todo~$A$. Assim, há ``mais'' !!{element}s em~$A$ do que
os infinitos inteiros positivos.

Como se provaria que um conjunto é !!{nonenumerable}? É preciso mostrar
que nenhuma tal função sobrejetora pode existir. Equivalentemente, é preciso
mostrar que os elementos de~$A$ não podem ser enumerados em uma lista
infinita unidirecional. A melhor maneira de fazer isso é mostrar que toda lista de
!!{element}s de~$A$ deve deixar de fora pelo menos um elemento; ou que nenhuma
função $f\colon \PosInt \to A$ pode ser !!{surjective}. Podemos fazer isso
usando o \emph{método diagonal} de Cantor. Dada uma lista de !!{element}s
de~$A$, digamos, $x_1$, $x_2$, \dots, construímos outro elemento de~$A$
que, por sua construção, não pode estar naquela lista.

Nosso primeiro exemplo é o conjunto~$\Bin^\omega$ de todas as sequências infinitas, sem lacunas,
de $0$s e $1$s.

\begin{thm}
\ollabel{thm:nonenum-bin-omega}
$\Bin^\omega$~é !!{nonenumerable}.
\end{thm}

\begin{proof}
Suponha, por absurdo, que $\Bin^\omega$ é
!!{enumerable}, isto é, suponha que há uma lista $s_{1}$, $s_{2}$,
$s_{3}$, $s_{4}$, \dots{} de todos os !!{element}s de~$\Bin^\omega$. Cada
um desses $s_i$ é em si uma sequência infinita de $0$s e~$1$s.
Chamemos o $j$-ésimo elemento da $i$-ésima sequência nessa lista de
$s_i(j)$. Então a $i$-ésima sequência~$s_i$ é
\[
s_i(1), s_i(2), s_i(3), \dots
\]

Podemos organizar essa lista, e os elementos de cada sequência $s_i$ nela,
em uma matriz:
\[
\begin{array}{c|c|c|c|c|c}
& 1 & 2 & 3 & 4 & \dots \\\hline
1 & \mathbf{s_{1}(1)} & s_{1}(2) & s_{1}(3) & s_1(4) & \dots \\\hline
2 & s_{2}(1)& \mathbf{s_{2}(2)} & s_2(3) & s_2(4) & \dots \\\hline
3 & s_{3}(1)& s_{3}(2) & \mathbf{s_3(3)} & s_3(4) & \dots \\\hline
4 & s_{4}(1)& s_{4}(2) & s_4(3) & \mathbf{s_4(4)} & \dots \\\hline
\vdots & \vdots & \vdots & \vdots & \vdots & \mathbf{\ddots}
\end{array}
\]
Os rótulos ao lado indicam o número da sequência na lista
$s_1$, $s_2$, \dots; os números no topo rotulam os !!{element}s
das sequências individuais. Por exemplo, $s_{1}(1)$ é um nome para
qualquer número, um $0$ ou um~$1$, que é o primeiro !!{element} na
sequência $s_{1}$, e assim por diante.

Agora construímos uma sequência infinita, $\overline{s}$, de $0$s e
$1$s que não pode estar nessa lista. A definição de
$\overline{s}$ dependerá da lista $s_1$, $s_2$, \dots. Qualquer
lista infinita de sequências infinitas de $0$s e $1$s dá origem a
uma sequência infinita~$\overline{s}$ que está garantidamente fora
da lista.

Para definir $\overline{s}$, especificamos todos os seus !!{element}s,
isto é, especificamos $\overline{s}(n)$ para todo $n \in \PosInt$. Fazemos isso
lendo a diagonal da matriz acima (daí o nome
``método diagonal'') e depois trocando cada $1$ por $0$ e cada
$0$ por~$1$. De maneira mais abstrata, definimos $\overline{s}(n)$ como $0$
ou $1$ conforme o $n$-ésimo !!{element} da diagonal,
$s_n(n)$, é $1$ ou $0$.
\[
\overline{s}(n) =
\begin{cases}
1 & \text{se $s_{n}(n) = 0$}\\
0 & \text{se $s_{n}(n) = 1$}.
\end{cases}
\]
Se você prefere fórmulas a definições por casos, também pode
definir $\overline{s}(n) = 1 - s_n(n)$.

Claramente $\overline{s}$ é uma sequência infinita de $0$s e
$1$s, já que é apenas a sequência espelhada da sequência de $0$s
e $1$s que aparecem na diagonal de nossa matriz. Assim $\overline{s}$
é !!a{element} de~$\Bin^\omega$. Mas não pode estar na lista $s_1$,
$s_2$, \dots{} Por quê?

Não pode ser a primeira sequência na lista, $s_1$, porque difere de
$s_1$ no primeiro !!{element}. Seja qual for $s_1(1)$, definimos
$\overline{s}(1)$ como o oposto. Não pode ser a segunda
sequência na lista, porque $\overline{s}$ difere de $s_2$ no segundo
elemento: se $s_2(2)$ é $0$, $\overline{s}(2)$ é $1$, e vice-versa. E assim por diante.

Mais precisamente: se $\overline{s}$ estivesse na lista, haveria
algum $k$ tal que $\overline{s} = s_{k}$. Duas sequências são idênticas
sse concordam em toda posição, isto é, para qualquer~$n$, $\overline{s}(n) =
s_{k}(n)$. Em particular, tomando $n = k$ como caso especial,
$\overline{s}(k) = s_{k}(k)$ teria de valer. $s_k(k)$ é $0$
ou~$1$. Se é $0$, então $\overline{s}(k)$ deve ser~$1$---é assim
que definimos $\overline{s}$. Mas se $s_k(k) = 1$, de novo por causa
da forma como definimos $\overline{s}$, $\overline{s}(k) = 0$. Em ambos
os casos $\overline{s}(k) \neq s_{k}(k)$.

Começamos supondo que há uma lista de !!{element}s de
$\Bin^\omega$, $s_1$, $s_2$, \dots{} A partir dessa lista construímos uma
sequência~$\overline{s}$ que provamos não poder estar na lista. Mas ela
definitivamente é uma sequência de $0$s e $1$s se todos os $s_i$ são
sequências de $0$s e $1$s, isto é, $\overline{s} \in
\Bin^\omega$. Isso mostra em particular que não pode haver lista de
\emph{todos} os !!{element}s de~$\Bin^\omega$, já que para qualquer tal lista
também poderíamos construir uma sequência~$\overline{s}$ garantidamente fora
da lista, de modo que a suposição de que há uma lista de todas as sequências
em~$\Bin^\omega$ leva a uma contradição.
\end{proof}

\begin{explain}
Esse método de prova é chamado de ``diagonalização'' porque usa a
diagonal da matriz para definir~$\overline{s}$. A diagonalização não
precisa envolver a presença de uma matriz: podemos mostrar que conjuntos não são
!!{enumerable} usando uma ideia semelhante mesmo quando não há matriz nem
diagonal propriamente dita.
\end{explain}

\begin{thm}
\ollabel{thm:nonenum-pownat}
$\Pow{\PosInt}$ não é !!{enumerable}.
\end{thm}

\begin{proof}
Procedemos da mesma forma, mostrando que para toda lista de subconjuntos
de~$\PosInt$ há um subconjunto de $\PosInt$ que não pode estar na lista.
Suponha que a seguinte é uma lista dada de subconjuntos de~$\PosInt$:
\[
Z_{1}, Z_{2}, Z_{3}, \dots
\]
Agora definimos um conjunto $\overline{Z}$ tal que para qualquer $n \in \PosInt$,
$n \in \overline{Z}$ sse $n \notin Z_{n}$:
\[
\overline{Z} = \Setabs{n \in \PosInt}{n \notin Z_n}
\]
$\overline{Z}$ é claramente um conjunto de inteiros positivos, já que por
suposição cada~$Z_n$ o é, e assim $\overline{Z} \in
\Pow{\PosInt}$. Mas $\overline{Z}$ não pode estar na lista. Para mostrar
isso, estabeleceremos que para cada $k \in \PosInt$, $\overline{Z} \neq
Z_k$.

Seja $k \in \PosInt$ arbitrário. Definimos $\overline{Z}$ de modo que
para qualquer $n \in \PosInt$, $n \in \overline{Z}$ sse $n \notin Z_n$.
Em particular, tomando $n=k$, $k \in \overline{Z}$ sse $k \notin Z_k$.
Mas isso mostra que $\overline{Z} \neq Z_k$, já que $k$ é !!a{element}
de um mas não do outro, e assim $\overline{Z}$ e $Z_k$ têm
!!{element}s diferentes. Como $k$ foi arbitrário, $\overline{Z}$ não está
na lista $Z_1$, $Z_2$, \dots
\end{proof}

\begin{explain}
A prova precedente não mencionou uma diagonal, mas você pode pensar nela
como envolvendo uma diagonal se a imaginar assim: imagine os
conjuntos $Z_1$, $Z_2$, \dots, escritos em uma matriz, onde cada
!!{element}~$j \in Z_i$ é listado na~$j$-ésima coluna. Digamos que os primeiros
quatro conjuntos nessa lista são $\{1,2,3,\dots\}$, $\{2, 4, 6, \dots\}$,
$\{1,2,5\}$ e $\{3,4,5,\dots\}$. Então a matriz começaria com
\[
\begin{array}{r@{}rrrrrrr}
  Z_1 = \{ & \mathbf{1}, & 2, & 3, & 4, & 5, & 6, & \dots\}\\
  Z_2 = \{ &  & \mathbf{2}, &  & 4, &  & 6, & \dots\}\\
  Z_3 = \{ & 1, & 2, &  &  & 5\phantom{,} &  & \}\\
  Z_4 = \{ &  &  & 3, & \mathbf{4}, & 5, & 6, & \dots\}\\
  \vdots & & & & & \ddots
\end{array}
\]
Então $\overline{Z}$ é o conjunto obtido percorrendo a diagonal,
deixando de fora quaisquer números que aparecem ao longo da diagonal e incluindo
aqueles $j$ onde a matriz tem uma lacuna na $j$-ésima linha/coluna. No
caso acima, deixaríamos de fora $1$ e $2$, incluiríamos~$3$, deixaríamos
de fora~$4$, etc.
\end{explain}

\begin{prob}
Mostre que $\Pow{\Nat}$ é !!{nonenumerable} por um argumento diagonal.
\end{prob}

\begin{prob}\label{sfr:siz:nen:prob:f-posint}
Mostre que o conjunto de funções $f \colon \PosInt \to \PosInt$ é
!!{nonenumerable} por um argumento diagonal explícito. Isto é, mostre que
se $f_1$, $f_2$, \dots, é uma lista de funções e cada $f_i\colon
\PosInt \to \PosInt$, então há alguma $\overline{f}\colon \PosInt \to
\PosInt$ fora dessa lista.
\end{prob}

\end{document}
